\section{Universal determinant completion and the intrinsic boundary atlas}
\label{sec:boundary-atlas-statement}

The residual-colon filtration of
Theorem~\ref{thm:global-residual-colon} is intrinsic at every
projection rank.  We first isolate the determinant-completion
mechanism in arbitrary dimension and symmetric degree.  We then
compute the projection-corank-two boundary on genuine generic
coefficient strata, including all pure-block rank drops, and finally
analyze projection coranks three and four.  The finite coefficient
matrices used below form a presentation atlas; the schemes \(W_j\),
their associated points, generic lengths, and graded multiplication
form the intrinsic boundary packet.

Keep the notation of
Theorem~\ref{thm:all-corank-presentation}.  Thus, at a point of
projection rank \(r\), put
\[
 H=\im(\mathcal K\to V),\qquad W=V/H,\qquad c=\dim W=4-r,
\]
choose a Schur chart
\[
 M=I_H\oplus T,\qquad d=\det T,
\]
and write
\begin{equation}\label{eq:atlas-residual-J}
 \mathcal I_{\widehat D_R}=dJ,\qquad
 J=I_6\!\left[
 A_H,\,
 B_H(1_H\otimes T),\,
 C_H\Sym^2T
 \right].
\end{equation}
The local nilradical is generated by the class of \(d\).

\begin{theorem}[Universal symmetric-power determinant completion]
\label{thm:universal-det-completion}
Let \(V\) be a complex vector space of dimension \(e\), let
\(r\ge1\), and let
\[
 \gamma:\Sym^rV\twoheadrightarrow S
\]
be a surjection, with \(s=\dim S\).  On \(\End(V)\) let \(M\)
be the generic endomorphism, put \(d=\det M\), and set
\[
 J_{\gamma,r}=I_s\!\left(\gamma\Sym^rM\right).
\]
Then
\begin{equation}\label{eq:universal-det-completion}
 d^{\kappa(e,r)}\in J_{\gamma,r},\qquad
 \kappa(e,r)=\binom{e+r-1}{r-1}.
\end{equation}
The exponent is sharp in the universal class: if
\(S=\Sym^rV\) and \(\gamma\) is an isomorphism, then
\(J_{\gamma,r}=(d^{\kappa(e,r)})\).
\end{theorem}

For \((e,r)=(4,2)\),
\(\kappa(4,2)=5\).  Thus the global depth bound below is a
specialization of a dimension-free determinant-completion principle,
rather than an isolated identity for the Hilbert function
\((1,4,6)\).

\begin{theorem}[Finite all-corank depth and multiplication]
\label{thm:finite-all-corank-depth}
For every \(R\in G_4^\circ\) and every projection-rank chart,
\begin{equation}\label{eq:d5-in-J}
 d^5\in J.
\end{equation}
Consequently
\[
 \widehat{\mathcal N}_R^{\,6}=0
\]
on the whole Grassmannian.  In particular
\[
 W_j=\varnothing\qquad(j\ge6),
\]
and the special fibre of the intrinsic normal-cone specialization has
at most five positive graded layers.

For \(i,j\ge1\), \(i+j\le5\), the multiplication map
\[
 \frac{\widehat{\mathcal N}^{\,i}}
      {\widehat{\mathcal N}^{\,i+1}}
 \otimes
 \frac{\widehat{\mathcal N}^{\,j}}
      {\widehat{\mathcal N}^{\,j+1}}
 \longrightarrow
 \frac{\widehat{\mathcal N}^{\,i+j}}
      {\widehat{\mathcal N}^{\,i+j+1}}
\]
is, under the identifications
\[
 \frac{\widehat{\mathcal N}^{\,j}}
      {\widehat{\mathcal N}^{\,j+1}}
 \simeq
 \mathcal L^{\otimes j}\otimes\mathcal O_{W_j},
\]
the tensor twist of the canonical quotient multiplication
\begin{equation}\label{eq:atlas-multiplication}
 \mathcal O_{W_i}\otimes_{\mathcal O_{\widehat\Delta_R}}
 \mathcal O_{W_j}
 \longrightarrow
 \mathcal O_{W_{i+j}}.
\end{equation}
In particular it is surjective.  Thus no additional multiplication
constant is hidden after the schemes \(W_j\) are known.
\end{theorem}

We next sharpen the corank-two analysis.  Assume \(A_H\) is
injective.  Put
\[
 E_H=S_R/\im A_H,\qquad
 \beta_H=\overline B_H:H\otimes W\longrightarrow E_H,\qquad
 \chi_H=\overline C_H:\Sym^2W\longrightarrow E_H.
\]
Here \(\dim H=\dim W=2\) and \(\dim E_H=3\).  Let
\[
 Q=\Pj(H\otimes W)_{\operatorname{rk}=1}
   \simeq\Pj^1\times\Pj^1
\]
be the Segre quadric.

\begin{definition}[Mixed-kernel orbit type]
\label{def:mixed-kernel-orbit-type}
The \emph{mixed-kernel orbit type} of \((R,H)\) is the rank
\(b=\rank\beta_H\), together with the incidence type of
\(\Pj(\ker\beta_H)\) with \(Q\).

For \(b=3\) this is a point off \(Q\) or on \(Q\).  For \(b=2\)
it is a secant line, a tangent line, an \(H\)-ruling
\(\Pj(h_0\otimes W)\), or a \(W\)-ruling
\(\Pj(H\otimes w_0)\) contained in \(Q\).  For \(b=1\), the one-dimensional image of
\(\beta_H^*:E_H^*\to(H\otimes W)^*\) is generated by a tensor of
rank two or rank one; equivalently the kernel plane cuts \(Q\) in a
smooth conic or in a pair of ruling lines.  The final type is
\(b=0\).
\end{definition}

These are exactly the nine
\(GL(H)\times GL(W)\times GL(E_H)\)-orbits for the mixed block.  The
first two are the mixed-regular and decomposable-kernel cases already
encountered in Section~\ref{sec:stratified-nilpotent}; the remaining
seven are higher-dimensional mixed kernels or failures of mixed
surjectivity.

For a homogeneous ideal \(K\subset
P=\C[a,b,c,d]\), write
\[
 H_{P/K}(n)=\dim_\C(P/K)_n.
\]
For each orbit, the \emph{generic coefficient stratum} is the dense
open obtained from a transverse stabilizer slice by deleting the
explicit denominator and leading-coefficient divisors in
Theorem~\ref{thm:generic-slice-certification}.  Equivalently it is
the locus on which the rational-function-field Groebner bases used
there retain their initial ideals.  The isolated integral matrices
listed in the proof are base points of those slices; they are not used
as a substitute for genericity.

\begin{theorem}[Corank-two mixed-kernel atlas]
\label{thm:corank-two-mixed-kernel-atlas}
Assume \(A_H\) is injective.  Then the first positive graded support
is determined without any mixed-regularity assumption:
\begin{equation}\label{eq:W1-all-Ainj}
 \mathcal O_{W_1}
 \simeq
 P/(\delta,g_0,\ldots,g_4),
 \qquad \delta=ad-bc,
\end{equation}
where, on \(P/(\delta)\) with \(T=uv^t\), the degree-four classes
span
\[
 h_H(u)\,\Sym^4(\C^2_v).
\]
Thus, if \(h_H\ne0\),
\[
 \Proj W_1=V(h_H)\times\Pj^1_v\subset Q,
\]
whereas \(h_H=0\) gives \(W_1=V(\delta)\).

On the generic coefficient stratum of each mixed-kernel orbit, the
remaining graded supports and the exact local nilpotency index are as
follows.
\begin{center}
\small
\begin{tabular}{@{}c l c l l@{}}
\toprule
\(b\)&mixed-kernel type&index&\(W_2\)&\(W_3\)\\
\midrule
3&point off \(Q\)&3&length \(1\)&empty\\
3&point on \(Q\)&3&length \(3\), \(H=(1,2)\)&empty\\
2&secant line&3&length \(8\), \(H=(1,4,3)\)&empty\\
2&tangent line&3&length \(8\), \(H=(1,4,3)\)&empty\\
2&\(H\)-ruling \(\Pj(h_0\otimes W)\)&4&length \(15\), \(H=(1,4,6,4)\)&length \(1\)\\
2&\(W\)-ruling \(\Pj(H\otimes w_0)\)&3&
 \(H(n)=n+1\) for \(n\ge4\)&empty\\
1&dual tensor rank \(2\)&4&
 \(H(n)=2n+2\) for \(n\ge3\)&length \(1\)\\
1&dual tensor rank \(1\)&4&
 \(H(n)=3n+3\) for \(n\ge3\)&length \(3\)\\
0&zero mixed map&4&\(V(\delta)\)&\(V(\delta)\)\\
\bottomrule
\end{tabular}
\end{center}
For \(b=0\), surjectivity of the original quotient multiplication
forces \(\chi_H\) to be an isomorphism and
\begin{equation}\label{eq:b0-delta3}
 J=(\delta^3);
\end{equation}
hence the last row is an exact identity, not only a generic one.

In the \(H\)-ruling and \(b\le1\) rows one has \(h_H=0\).  On the
\(W\)-ruling row the contact quartic is nonzero on the transverse open
and has the root forced by \(w_0\).
For \(b=1\), the projectivization of \(W_2\) has degree two in the
dual-rank-two case and degree three in the dual-rank-one case.  The
length-one and length-three third layers are respectively the reduced
rank-two vertex and the length-three vertex thickening already
appearing on the decomposable-kernel wall.
\end{theorem}

The complement of the generic coefficient stratum is itself part of
the atlas, rather than an omitted exceptional set.  We make this
finite refinement explicit.

\begin{definition}[Macaulay valuation atlas]
\label{def:macaulay-atlas}
Let \(P_c=\C[t_{ij}:1\le i,j\le c]\), \(d=\det T\), and let \(J\)
be the residual ideal on a fixed projection-rank chart.  For
\(1\le k\le5\), let
\[
 \mathsf M_k(J)
\]
be the coefficient matrix whose columns are all products
\[
 m\,q,\qquad
 q\in\mathcal G(J),\qquad
 \deg m+\deg q=ck,
\]
written in the monomial basis of \((P_c)_{ck}\), where
\(\mathcal G(J)\) is the displayed set of maximal minors in
\eqref{eq:atlas-residual-J}.  Let \(v_k\) be the coefficient vector
of \(d^k\).  Define
\[
 \varepsilon_k(J)=
 \rank[\mathsf M_k(J)\mid v_k]-\rank\mathsf M_k(J)\in\{0,1\}.
\]
Then
\[
 \varepsilon_k(J)=0\quad\Longleftrightarrow\quad d^k\in J.
\]
The locally closed loci obtained by fixing the ranks of the matrices
\(\mathsf M_k\), their augmented matrices, and the analogous matrices
for the kernels
\[
 K_j=(J:d^{\,j-1}),\qquad1\le j\le5,
\]
are called the \emph{Macaulay valuation atlas}.  All of their
equations are minors of explicit finite coefficient matrices.
\end{definition}

\begin{theorem}[Bounded principal-colon primary stratification]
\label{thm:bounded-principal-colon-stratification}
Let \(S\) be noetherian, let \(A\) be a finitely presented
\(S\)-algebra which is flat over \(S\), let \(J\subset A\) be a
finitely generated ideal, and let \(f\in A\).  Assume that
\(f^N\in J\) for some \(N\ge1\).  Put
\[
 B=A/J,\qquad
 K_q=(J:f^q),\qquad
 M_q=A/(f,K_q),\qquad 0\le q<N.
\]
Then \(S\) admits a finite locally closed stratification
\(S=\coprod_\alpha S_\alpha\) with the following properties.

For every \(\alpha\), all the assertions below are preserved by an
arbitrary base change \(T\to S_\alpha\).
\begin{enumerate}
\item Formation of \(J\) and of every \(K_q\) commutes with base
change.  The modules
\[
 B,\qquad C_q=A/(J,f^q),\qquad
 I_q=f^qB,\qquad M_q
\]
are flat over \(S_\alpha\).
\item For each \(q\) there are finitely many closed relative support
packets \(Z_{q,\nu}\subset\Spec A_{S_\alpha}\), flat after a further
refinement, such that for every geometric point \(\bar s\to S_\alpha\)
the associated points of \((M_q)_{\bar s}\) are exactly the generic
points of the irreducible components of the nonempty
\((Z_{q,\nu})_{\bar s}\).  Geometric splitting of a packet is allowed.
Dimensions, incidence, and minimal-versus-embedded status are constant
on the stratum.
\item If \(p\) is one of these geometric associated points, the
canonical embedded multiplicity
\[
 \mu_p(M_q)=
 \operatorname{length}_{\OO_{\Spec A_{\bar s},p}}
 H^0_{p\OO_{\Spec A_{\bar s},p}}\bigl((M_q)_p\bigr)
\]
is constant along the corresponding packet.
\item Any prescribed finite collection of maps between the \(M_q\)
may be included in the same refinement so that the relevant source,
target, image and cokernel modules are flat; in particular their
fibre ranks commute with base change.
\end{enumerate}
Thus a bounded principal power converts the a priori infinite
primary problem into a finite, base-change-compatible invariant.
\end{theorem}

\begin{corollary}[Universal primary finiteness for symmetric-power quotients]
\label{cor:universal-primary-finiteness}
Fix integers \(e,r\ge1\) and \(s\le\dim\Sym^rV\).  Let \(\mathcal U\)
be the open parameter space of surjections
\[
 \gamma:\Sym^rV\twoheadrightarrow S,\qquad \dim S=s,
\]
and let \(M\) be the universal endomorphism of \(V\).  On
\(\mathcal U\times\End(V)\) put
\[
 J_{\mathrm{univ}}=
 I_s\!\left(\gamma\Sym^rM\right),\qquad d=\det M,
 \qquad N=\binom{e+r-1}{r-1}.
\]
Then, after a finite locally closed stratification of \(\mathcal U\),
formation of every residual colon
\[
 (J_{\mathrm{univ}}:d^q),\qquad 0\le q<N,
\]
commutes with arbitrary base change.  The corresponding modules
\[
 \mathcal O/(d,(J_{\mathrm{univ}}:d^q))
\]
admit finitely many geometric associated-support packets with constant
minimal/embedded incidence, constant embedded multiplicities, constant
Hilbert polynomials, and base-change-compatible multiplication ranks.
Thus the universal determinant exponent controls not merely nilpotence
depth but the finite primary complexity of the entire symmetric-power
determinantal family.
\end{corollary}

\begin{proof}
Theorem~\ref{thm:universal-det-completion} gives
\(d^N\in J_{\mathrm{univ}}\) identically on the universal
surjection space.  Apply
Theorem~\ref{thm:bounded-principal-colon-stratification} to the
finitely presented flat coefficient algebra on
\(\mathcal U\times\End(V)\), with \(f=d\) and
\(J=J_{\mathrm{univ}}\).  All stated conclusions are exactly the
four clauses of that theorem, taken simultaneously for the finite list
\(0\le q<N\).
\end{proof}

The theorem is independent of the special \((1,4,6)\) model.  In the
present paper the universal determinant-completion theorem supplies
\(N=5\) on every projection-rank chart.  Consequently the residual
colon algebra, its associated-point packets, and all graded
multiplication maps admit one finite primary stratification before
any orbit-specific computation is made.
Because of Theorem~\ref{thm:finite-all-corank-depth}, no matrix with
\(k>5\) occurs.  We refine these rank strata simultaneously by the
flattening strata of the finite modules
\[
 \mathcal M_j=P_c/\bigl((d)+K_j\bigr),\qquad 1\le j\le5,
\]
by the flattening strata of the cokernels of
\eqref{eq:atlas-multiplication}, and by the associated-point
stratification in the following proposition.

\begin{proposition}[Finite intrinsic primary-signature refinement]
\label{prop:associated-prime-refinement}
On every finite-type coefficient chart for
\((A_H,B_H,C_H)\), the presentation atlas admits a finite locally
closed refinement such that the following statements hold for every
geometric fibre.
\begin{enumerate}
\item Formation of each residual colon
\(K_j=(J:d^{j-1})\), \(1\le j\le5\), commutes with base change,
and
\[
 \mathcal M_j=P_c/\bigl((d)+K_j\bigr)
\]
is flat over the stratum with constant Hilbert polynomial.
\item There are finitely many relative closed supports
\(Z_{j,\nu}\) whose nonempty fibrewise generic points are exactly
\(\Ass((\mathcal M_j)_s)\).  Their dimensions,
minimal-versus-embedded status, incidence relations, and generic
lengths are constant on the stratum.
\item The cokernels of the graded multiplication maps are flat, so
their ranks are constant and commute with base change.
\end{enumerate}
Consequently the radicals, associated supports, generic
multiplicities, nilpotency index, Hilbert polynomials, and
multiplication ranks of all positive graded layers have constant
intrinsic primary signature on each final stratum.

The refinement is not asserted to be canonical as a collection of
coordinate strata.  The primary signature is intrinsic; different
Schur charts, generating sets, or monomial orders admit a common
finite refinement and compute the same signature.
\end{proposition}

This refinement is used below for rank drops.  The proof in
Section~\ref{sec:relative-primary-specialization} is the application of
Theorem~\ref{thm:bounded-principal-colon-stratification}: it flattens
both \(P_S/J\) and \(P_S/(J,d^{j-1})\), deduces flatness of the image
of multiplication by \(d^{j-1}\), and only then invokes relative
assassins.  Thus coefficient collisions are
recorded by actual specialization of the intrinsic graded algebra,
not by relabelling a primary decomposition.

There is also a useful degree bound that detects many strata before
any Macaulay matrix is formed.  Put
\[
 a=\rank A_H,\qquad q=6-a,
\]
and, after quotienting by \(\im A_H\),
\[
 b=\rank\overline B_H.
\]
Every nonzero residual \(q\)-minor uses at most \(b\) linear mixed
columns and therefore has transverse degree at least
\begin{equation}\label{eq:minimal-transverse-degree}
 \ell(a,b)=2q-b.
\end{equation}

\begin{proposition}[Rank-drop depth bounds at corank two]
\label{prop:rankdrop-coranktwo-depth}
At projection corank two,
\begin{equation}\label{eq:coranktwo-lower-bound}
 \min\{k:d^k\in J\}\ge
 \kappa_0(a,b):=
 q-\left\lfloor\frac b2\right\rfloor.
\end{equation}
Together with \(d^5\in J\), this leaves only the following finite
possibilities when \(A_H\) drops rank:
\[
\begin{array}{c|c|c}
a&q&\text{possible }(b;\,k)\\ \hline
2&4&(4;\,2,3,4,5),\ (3,2;\,3,4,5),\ (1;\,4,5),\\
1&5&(4;\,3,4,5),\ (3,2;\,4,5),\\
0&6&(4;\,4,5),\ (3;\,5).
\end{array}
\]
Here \(k=\min\{j:d^j\in J\}\), so the nilpotency index is \(k+1\).
The omitted values of \(b\) cannot occur under surjectivity of the
original quadratic quotient, because the quadratic block contributes
at most three independent target directions.
Each comma-separated choice of \(k\) is cut out exactly by
\(\varepsilon_k(J)=0\) and
\(\varepsilon_{k-1}(J)=1\).
In particular the extreme rank-drop type \((a,b)=(0,3)\) has exact
nilpotency index six.
\end{proposition}

At projection corank three a stronger representation-theoretic
exclusion reduces the atlas to one binary choice.

\begin{theorem}[Projection corank three]
\label{thm:corank-three-depth}
At every projection-corank-three point of \(\widehat D_R\),
including all rank drops of \(A_H\) and of the mixed block,
\begin{equation}\label{eq:corank3-no-d3}
 d^3\notin J,\qquad d^5\in J.
\end{equation}
Hence the nilpotency index is either five or six.  It is five exactly
when
\[
 \varepsilon_4(J)=0,
\]
and six exactly when \(\varepsilon_4(J)=1\).
Thus the corank-three locus has two intrinsic determinant-membership
strata.  Any Macaulay presentation atlas computes these strata, but
the chosen coefficient matrices are not themselves declared
canonical.
\end{theorem}

At the deepest Schubert stratum the remaining condition is identified
geometrically.

\begin{theorem}[Projection corank four and the Jacobian quartic]
\label{thm:corank-four-jacobian}
Let \(R\in G_4^\circ\).  At a projection-corank-four point,
\(H=0\), \(T\in\operatorname{Mat}_{4\times4}\), and
\[
 J=I_6(\gamma_R\Sym^2T).
\]
Then
\begin{equation}\label{eq:corank4-depth}
 d^3\notin J,\qquad d^4\in J,\qquad d=\det T.
\end{equation}
Consequently the local nilradical has exact index five.

More precisely, as a \(GL(V)\)-module,
\begin{equation}\label{eq:wedge6-sym2-decomposition}
 \bigwedge\nolimits^6\Sym^2V
 \simeq
 \mathbb S_{(5,4,2,1)}V
 \oplus
 \mathbb S_{(4,4,4,0)}V.
\end{equation}
The absence of \((\det V)^3\) in
\eqref{eq:wedge6-sym2-decomposition} gives the first
nonmembership in \eqref{eq:corank4-depth}.  Under the canonical
duality
\[
 \bigwedge\nolimits^6\Sym^2V
 \simeq
 (\det V)^5\otimes
 \left(\bigwedge\nolimits^4\Sym^2V\right)^*,
\]
the second summand in \eqref{eq:wedge6-sym2-decomposition} is dual,
up to determinant twist, to
\[
 \det V\otimes\Sym^4V,
\]
and the corresponding covariant of \(\bigwedge^4R\) is the Jacobian
quartic \(F_R\).  Since \(F_R\ne0\) on \(G_4^\circ\), the determinant
line \((\det V)^4\) occurs in the degree-sixteen part of \(J\), which
is exactly the second membership in
\eqref{eq:corank4-depth}.
\end{theorem}

Theorems~\ref{thm:finite-all-corank-depth},
\ref{thm:corank-two-mixed-kernel-atlas},
\ref{thm:corank-three-depth}, and
\ref{thm:corank-four-jacobian}, together with the generic primary
signature and specialization theorems of
Section~\ref{sec:relative-primary-specialization}, form the boundary
classification used in this revision.  The residual-colon formulas
are coordinate-free identities; the presentation matrices are finite
computational devices.  At corank two the generic supports and their
primary corrections are geometrically identified, including pure-block
rank drops.  At higher corank the exact depth statements remain
intrinsic and are coupled to the finite primary-signature refinement.
